\documentclass[11pt]{article}

% ----- Encoding & fonts -----
\usepackage[utf8]{inputenc}
\usepackage[T1]{fontenc}
\usepackage{mathptmx}   % scalable Times text + math (Type 1)
\usepackage{courier}    % scalable typewriter (for \texttt network names)

% ----- Layout -----
\usepackage[margin=1in]{geometry}
\usepackage{parskip}
\usepackage{microtype}

% ----- Math -----
\usepackage{amsmath}
\usepackage{amssymb}

% ----- Tables -----
\usepackage{booktabs}
\usepackage{array}
\usepackage{longtable}

% ----- Lists & links -----
\usepackage{enumitem}
\usepackage[hidelinks]{hyperref}
\hypersetup{
  pdftitle={Conservation-by-Construction GNN Emulator of Silicon Burning: Consolidated Architecture Specification},
  pdfsubject={Synthesis of three architecture design runs and three adversarial referee audits},
  pdfkeywords={graph neural network, silicon burning, nucleosynthesis, conservation, emulator}
}

% ----- Convenience macros -----
\DeclareMathOperator{\asinh}{asinh}
\newcommand{\msn}[1]{\texttt{mesa\_#1}}
\newcommand{\Ye}{Y_{e}}
\newcommand{\dYe}{\Delta Y_{e}}
\newcommand{\enuc}{e_{\mathrm{nuc}}}
\newcommand{\dd}{\mathrm{d}}
\newcommand{\Tr}{^{\mathsf{T}}}

\title{Conservation-by-Construction GNN Emulator of Silicon Burning\\[4pt]
\large Consolidated Architecture Specification (Components A--D) with Adversarial Verification}
\author{}
\date{June 2026}

\begin{document}
\maketitle

{\small\itshape
Synthesis of six research artifacts --- three architecture design runs (Component A: graph backbone and message passing; Component B: output heads, conservation map, equilibrium masking, energy/neutrino; Components C+D: variable-$\Delta t$ temporal head, rollout stability, and training recipe) and their three adversarial referee audits, June 2026. All audit corrections are folded directly into the text below; figures, citations, and numerical bands reflect the post-verification state. The design rests on a separately established foundation --- the project's scientific case and the prediction-target/$\Ye$-accuracy decision --- which is summarized only as far as needed to make the architecture self-contained.\par}

\begin{abstract}
This report consolidates the full architectural design of a conservation-by-construction graph neural network (GNN) emulator for the silicon-burning / core-collapse-supernova (CCSN) progenitor regime, together with three rounds of adversarial verification of that design. The emulator targets the same problem as Grichener et al.'s 2025 ``Nuclear Neural Network'' (NNN) --- fast surrogate evaluation of MESA's softwired nuclear burning (\msn{80}, \msn{151}) on its public Zenodo data --- but replaces the NNN's dense fully-connected architecture with a heterogeneous graph operating on the bipartite isotope$\leftrightarrow$reaction graph, and adds the property the NNN and the closest competitor (NuGNN) both lack: hard baryon-number and charge conservation built into the architecture, with the electron fraction $\Ye$ free to evolve correctly through the weak interactions that drive it.

The architecture decomposes into three composable pieces. \textbf{Component A} is an encode--process--decode backbone on the bipartite graph with per-reaction-type message functions, GATv2 attention, and a feature set that adds the four physics channels silicon burning needs and NuGNN omits (electron screening, temperature-dependent partition functions, $\rho\Ye$ conditioning for tabulated weak rates, and reverse-rate/detailed-balance linkage), with size transfer carried by a continuous $(Z,N)$ embedding rather than isotope-index one-hots. \textbf{Component B} makes Target A --- a signed net per-reaction flux $\varphi$ passed through a fixed stoichiometric matrix, $\dd Y = \nu\varphi$ --- the primary head, so that baryon and charge conservation hold to machine precision for any $\varphi$ (the Sturm--Wexler flux-target / fixed-linear-layer result); a hybrid equilibrium mask (a cheap Guidry partial-equilibrium prior refined by a learned $L_0$ hard-concrete gate, acting on the latent flux and never on weak reactions) keeps the head from resolving the near-zero net flux of an equilibrated forward/reverse pair; and the energy release is computed from the same net fluxes, $\enuc = \sum_j Q_j\varphi_j$, which is automatically consistent with the composition update. \textbf{Components C and D} ship $\Delta t$ as a global conditioning feature on a stiffness-aware logarithmic grid (one model spanning $10^{-6}$--$10^{2}$~s, versus the NNN's nine), keep all conservation in the decode step rather than inside any latent dynamics, govern the rollout with a timescale-based rescaled update that refuses to trust error-dominated short steps, and train with pushforward/unrolled losses, a worst-case tail term, up-weighting of the Fe-peak weak-rate nuclei, and a deep-ensemble out-of-distribution gate that hands hard zones back to the solver.

All three citation bases survived hostile, line-by-line verification with no hallucinated load-bearing references; the corrections required were matters of attribution, labeling, and internal consistency, and are incorporated silently here. What the audits leave standing are a small number of genuine, load-bearing uncertainties that no literature can resolve and that define the project's Phase-0 measurement program: whether size transfer to chemically new neutron-rich isotopes holds, whether the per-step $\Ye$ residual accumulates systematically or as a random walk (which moves the pass/fail gate by two orders of magnitude), whether error distributions measured on quasi-random Sobol compositions transfer to the thin manifold of real MESA trajectories, whether the trajectory-varying equilibrium mask is stable enough for rollout training, and whether the asinh latent and the out-of-regime rollout governor behave at silicon-burning dynamic range. The overall verdict across the three audits is consistent: the architecture is sound and the niche is open as of mid-June 2026; commit to the design conditionally, with the measurements above as gates and with explicit thresholds at which to switch the target, the processor, the mask, or the temporal head rather than persist.
\end{abstract}

\section{Inherited foundation (the established basis the architecture builds on)}

The architecture is downstream of two settled decisions that this report does not re-derive but states for completeness.

The \textbf{scientific case} is that the pre-collapse core structure --- and hence the neutron-star-versus-black-hole outcome, explosion energy, and yields --- is controlled by the electron fraction $\Ye \equiv \sum_i Z_i X_i / A_i = \sum_i Z_i Y_i$, because the effective Chandrasekhar mass scales as $M_{\mathrm{ch}}^{\mathrm{eff}} \propto \Ye^{2}$. Accurate $\Ye$ during oxygen and silicon burning requires weak interactions (electron captures and $\beta$-decays) that are absent from the small $\alpha$-chain networks used in most stellar-evolution runs; Farmer et al.\ 2016 find $\approx$30\% variation in central $\Ye$ and shell locations across network choices and a $\approx$127-isotope floor for $\approx$10\% convergence. The reason large networks are rarely run in full is stiffness, which forces implicit integration with repeated Jacobian inversions. The emulator therefore competes not only against the full solver but against the QSE-reduced networks (Hix et al.\ 2007, roughly an order of magnitude faster with under a third the variables) and the small-net-plus-post-processing workflow.

The \textbf{prediction-target decision} is that conservation by construction is the right differentiator but must be realized so that it does not walk into the quasi-statistical-equilibrium (QSE) catastrophic-cancellation regime that defines silicon burning. Above $\sim$3~GK the fast capture/photodisintegration pairs balance and a species' net change is the small difference of two enormous gross fluxes; predicting non-negative forward and reverse fluxes and subtracting them is exactly the operation the QSE literature has warned against since Hix \& Thielemann 1996. The resolution, validated cross-domain in atmospheric photochemistry (Sturm \& Wexler) and combustion (CRNN, D\"oppel--Votsmeier), is that \textbf{the conserved quantity must be the network's native linear output, with any nonlinear transform kept internal/latent}. Two parameterizations satisfy this: predict the signed net flux directly and map it through the fixed stoichiometric matrix (Target A), or predict $\dd X$ directly and project onto the conservation null-space in a linear space (Target B). Target A is the lean; Target B is the benchmark and fallback.

The \textbf{accuracy gate} is set by the weak-interaction physics that drives $\Ye$. The defensible nuclear-physics floor, anchored on the FFN$\to$LMP rate-set shift and led by the Fe-peak electron-capture/$\beta$-decay nuclei (notably around ${}^{55}\mathrm{Co}$), is $\dYe \approx 5\times10^{-3}$ to $1.5\times10^{-2}$ end-to-end per trajectory. Under the conservative systematic-(linear)-accumulation model appropriate to an autoregressive emulator, this implies a per-step kill-test threshold of $|\dYe| \lesssim 3\times10^{-6}$ at a trajectory length $N \approx 1.6\times10^{3}$. This $3\times10^{-6}$ figure is the \textbf{operative} per-step gate throughout the architecture below; the alternatives implied by a longer trajectory ($5\times10^{-7}$ systematic at $N = 10^{4}$) or by zero-mean accumulation ($\approx 5\times10^{-5}$ random-walk at $N = 10^{4}$) are not adopted until the accumulation model is measured (Section~4.4).

The \textbf{regime box} is the Grichener Sobol domain: $T \approx 1.6$--$7.9$~GK ($10^{9.2}$--$10^{9.9}$~K), $\rho \approx 10^{7}$--$10^{9}$~g\,cm$^{-3}$, $0.45 < \Ye < 0.5$; the silicon-burning core of this box (Hix \& Thielemann 1996) sits at $T \approx 3.5$--$5$~GK, $\rho \approx 10^{7}$--$10^{10}$~g\,cm$^{-3}$, $\Ye \approx 0.498$--$0.46$. Three nuclear pools coexist here --- a light-particle pool ($n$, $p$, $\alpha$, light nuclei), a silicon group ($A \approx 24$--$45$), and an iron-peak group ($A \approx 50+$) --- separated by the $Z = N = 20$ shell closures; the net flow between the silicon and iron-peak groups is carried predominantly by the single bottleneck reaction ${}^{45}\mathrm{Sc}(p,\gamma){}^{46}\mathrm{Ti}$ ($\approx$75\% of the flow at the proton-rich $Z = 21$ edge), which is therefore the reaction whose flux the Phase-0 kill-test instruments.

The two reference points are the \textbf{baseline} --- Grichener et al.\ 2025 (ApJS 279, 49), a dense MLP emulator of MESA's bbq burning with a log-softmax composition head that enforces $\sum X_i = 1$ only (not baryon or charge conservation), a \emph{single shared network} emitting composition plus normalized energy and neutrino terms (the two-separate-network design was tried and abandoned), nine per-timestep models with runtime interpolation, large per-isotope errors reported as means, rollout error accumulation, and no MESA coupling --- and the \textbf{closest competitor}, NuGNN (Kim et al.\ 2026, arXiv:2606.04491, preprint), the first GNN on the bipartite isotope/reaction graph, demonstrated on a 690-isotope X-ray-burst network, which validates the graph inductive bias (signed-log MAE 0.037 versus 0.11 for a Res-U-Net and 0.56 for an FNN; in-solver substitution reproduces final abundances where both baselines fail) but predicts $\dd X$ without hard conservation and reports abandoning a conservation loss/activation because its signed-log target preprocessing made the inverse transform unstable.

\section{Component A --- graph backbone and heterogeneous message passing}

\subsection{Graph and topology}

The graph is exported from the pynucastro NetworkX representation as a heterogeneous bipartite structure: isotope nodes $I$ and reaction nodes $R$, with three edge types --- $I\!\to\!R$ (reactant incidence), $R\!\to\!I$ (signed stoichiometric incidence, carrying the sign of $\nu_{ij}$), and $I\!\to\!I$ (intra-reaction coupling). This is NuGNN's topology, and the bipartite reaction node \emph{is} the natural hypergraph encoding: the reaction node materializes the hyperedge and the incident edges carry the stoichiometry. Explicit hypergraph convolution (ChemHGNN, Rxn-Hypergraph) is rejected because nuclear reactions in the REACLIB chapters have arity $\le 3$ (triple-$\alpha$ is the extreme), so the incidence-normalization and hyperedge-degree machinery buys nothing and complicates per-reaction-type weight sharing. The fixed stoichiometric matrix $\nu$ is retained externally for Component B's conservation map. (Switch rule: adopt hypergraph convolution only if the network contains a meaningful population of $\ge 4$-species reactions \emph{and} an ablation shows reaction-node message passing loses $>0.01$ signed-log MAE on them --- unlikely to trip for softwired \msn{80/151/204}.)

\subsection{Feature design}

This is the most consequential area, because the closest template (NuGNN) was tuned for a proton-rich, rp/$\alpha$p-process X-ray-burst regime, whereas silicon burning is neutron-rich, QSE-dominated, and weak-rate-driven. The design feeds \textbf{pre-evaluated signed-log rates at the current $(T, \rho\Ye)$} rather than raw REACLIB coefficients, so the network is handed the physically meaningful quantity instead of being forced to relearn $\lambda(T)$ across 1.6--7.9~GK (this matches NuGNN's own practice; pynucastro evaluates the rates analytically and cheaply). It then adds the four channels NuGNN omits that the $\Ye$ regime requires:

\begin{enumerate}[leftmargin=*]
  \item \textbf{Electron screening factors.} At $\rho \approx 10^{7}$--$10^{9}$~g\,cm$^{-3}$ charged-particle rates are screening-enhanced, and screening shifts the QSE distribution itself (Hix \& Thielemann 1996). pynucastro's generated networks implement weak/intermediate/strong screening (Graboske 1973; Alastuey \& Jancovici 1978; Itoh 1979).
  \item \textbf{Temperature-dependent partition functions $G(T)$}, required for correct reverse rates above 3~GK (Rauscher \& Thielemann 2000; pynucastro supplies Rauscher 2003 tables with FRDM defaults).
  \item \textbf{Explicit $\rho\Ye$ conditioning on weak-reaction nodes.} The electron-capture/$\beta$-decay rates that drive $\Ye$ are tabulated in $(T, \rho\Ye)$, not as REACLIB fits, and the final $\Ye$ is set during silicon shell burning before the core's final Kelvin--Helmholtz contraction (Heger et al.\ 2001, ApJ 560:307; LMP rates raise central $\Ye$ by $\approx$0.01--0.015 over the older FFN set). $\rho\Ye$ must therefore be an explicit feature on weak nodes, masked off on non-weak nodes.
  \item \textbf{Reverse-rate / detailed-balance linkage.} Above $\sim$3~GK forward/reverse pairs near-cancel, so for Target A the network must see the pre-computed net (or both members of each pair) to keep the QSE near-cancellation internal rather than forming a difference of large learned numbers (Rauscher \& Thielemann 2000 for the detailed-balance construction; the QSE temperature threshold itself is Hix \& Thielemann 1999).
\end{enumerate}

The implementable feature vectors are: isotope node $x_I \in \mathbb{R}^{24}$ --- 16 physics/state channels [$\asinh X$, signed-log net flux, isotope-mapped signed-log rate aggregate, 1p/2p/1n/2n separation energies, mass-excess (binding energy per $A$), $\log G(T)$, $\log T$, $\log\rho$, $\log\Delta t$, $\Ye$ broadcast, screening aggregate, $Z/A$, $N/A$] concatenated with an 8-d $(Z,N)$ embedding; reaction node $x_R \in \mathbb{R}^{16}$ --- 8 physics channels [signed-log evaluated rate, $Q$-value, $\log T$, $\log\rho$, $\log\Delta t$, $\log(\rho\Ye)$ on weak reactions else masked, screening factor on charged-particle reactions, reverse-rate value-or-flag] concatenated with an 8-d reaction-type embedding; and a global $g \in \mathbb{R}^{3} = [\log T, \log\rho, \log\Delta t]$ with $\Ye$ broadcast to nodes, injected by FiLM into each processor block. For Target A, the reverse-rate pair value is carried explicitly on the reaction node so the near-cancellation lives as an internal latent quantity; for Target B it is less load-bearing because $\dd X$ is the directly supervised linear output. Each added channel must improve $\Ye$ error or signed-log MAE by $\ge 0.005$ absolute in Phase-0 ablation to be retained.

The \textbf{reaction-type taxonomy} extends NuGNN's seven XRB-centric types --- $(p,\gamma)$, $(\gamma,p)$, $(\alpha,\gamma)$, $(\gamma,\alpha)$, $(p,\alpha)$, $(\alpha,p)$, $\beta^{+}$ --- with the channels the Si/neutron-rich regime adds: \textbf{electron-capture, $(n,\gamma)/(\gamma,n)$, and three-body (triple-$\alpha$)}. Omitting these would leave the $\Ye$-driving and $\alpha$-network reactions untyped.

\subsection{Message-passing processor and depth}

The backbone is an encode--process--decode design (Battaglia 2018; GNS, Sanchez-Gonzalez 2020; MeshGraphNets, Pfaff 2020): a two-layer per-node-type encoder to latent width \textbf{$h = 128$} (wider than NuGNN's 64 to carry the extra physics channels and serve both targets --- revisited in Phase-0, as over-width risks overfitting the narrow $\Ye$ band), a processor of message-passing blocks with per-reaction-type message and update functions, and a decoder deferred to Component B. Attention uses \textbf{GATv2} (strictly more expressive than GAT for heterogeneous, feature-conditioned edges), dropped to mean/sum aggregation if it gives $<0.003$ MAE improvement.

The processor's block-sharing is the one place where the backbone choice is driven by a downstream consideration, and the verified position is to make the \textbf{untied multi-block stack the default} --- it is the only configuration in which NuGNN's accuracy floor was actually achieved --- and treat the \textbf{weight-shared recurrent processor as the contingent option}, adopted only if Phase-0 shows it stays within 0.01 absolute signed-log MAE of untied blocks at equal step count. Compatibility with a Component-C latent-ODE head (which a recurrent processor would enable) is a tiebreaker, not a driver.

Depth is set by the graph radius: $K \approx \lceil\text{graph radius}\rceil + 2$, starting from $K \approx 6$ for \msn{80} and growing as the network grows, with the radius measured in Phase-0 rather than assumed (the QSE clusters are tightly internally coupled and linked by few bridging reactions, so the radius is expected to be small). Because the topology includes direct $I\!\to\!I$ edges (1 hop), the effective reach is larger than a reaction-mediated ``$K$ steps reach $\sim K/2$ reaction-links'' estimate would suggest; the depth-reach reasoning accounts for those 1-hop edges (or the $I\!\to\!I$ edges are dropped if Phase-0 shows they do not help). \textbf{Depth is gated on the validation-MAE plateau} (freeze $K$ when $K\!\to\!K{+}1$ improves MAE by $<0.002$) and the measured radius --- not on a global Dirichlet-energy over-smoothing alarm, which is confounded in exactly this regime: within-group feature homogenization is the physical signature of QSE equilibration, not pathology, so a global Dirichlet collapse is expected and physical. If a smoothing diagnostic is wanted at all, it is computed per-QSE-group or restricted to non-equilibrated reaction subsets.

\subsection{Symmetry and size transfer}

The graph lives on the $(N,Z)$ chart of nuclides, not in 3D space, so there is no rotational or translational symmetry to exploit; $(N,Z)$ are physically meaningful labels, and only permutation equivariance over nodes is required, which plain heterogeneous message passing already provides. E$(n)$/SE(3)-equivariant machinery is rejected as overhead that could impose false isotropy. The weak exploitable structure --- local smoothness of nuclear properties on the $(N,Z)$ lattice --- is captured by the learned embeddings and separation-energy features. (Switch rule: adopt lattice-translation-equivariant layers only if Phase-0 shows $(N,Z)\!\to\!(N{+}1,Z{+}1)$ weight-tying improves new-isotope transfer by $\ge 0.01$ MAE.)

Size transfer is carried by \textbf{shared node-type and reaction-type embeddings plus a learned continuous $(Z,N)$ embedding MLP} ($\mathbb{R}^{2}\!\to\!\mathbb{R}^{8}$), not an isotope-index lookup, so a never-seen isotope still receives a smooth embedding by interpolation, and by feeding \emph{physical} features (separation energies, mass excess, partition functions, screening) that are defined for new isotopes. The epistemic discipline here is the report's strongest move and survives audit cleanly: \textbf{count transfer is established} (GNS generalizes to $\ge 1$ order of magnitude more particles at test; MeshGraphNets scales to more complex state spaces), but \textbf{transfer to chemically new neutron-rich isotopes is a hypothesis}, because the \msn{151}$\,\setminus\,$\msn{80} species carry \emph{new weak physics} exactly where \msn{80} fails --- they are not ``more of the same node.'' The hypothesis is falsified if zero-shot $\Ye$ error on \msn{151} exceeds the \msn{80}-internal $\Ye$ error by more than $\approx 2\times$ (degrading the very quantity the larger network exists to fix); if falsified, the model is reported as size-\emph{adaptable} via few-shot fine-tuning, not size-\emph{transferable} zero-shot. This is the single biggest risk in the backbone.

\subsection{Backbone normalization and the decode contract}

Signed-log is used for fluxes, rates, and $\dd X$, with NuGNN's constant-shift trick $\mathrm{sign}(\cdot)\,(C + \log_{10}|\cdot/X|)$, $C \approx 17$, to preserve magnitude ordering; per-channel standardization is fit on Phase-0 data; inference is float32 with a float64 cast for the $X_{t+\Delta t} = X_t + \Delta X$ update (NuGNN-verified safe). NuGNN's learned flux preprocessor is retained and is arguably more important here than for NuGNN, since it directly refines the signed net flux $\varphi$ that Target A predicts. The backbone is target-agnostic and exposes per-isotope latents $h_I$ and per-reaction latents $h_R$; the contract with Component B is that Target A reads $\varphi$ from $h_R$ and Target B reads $\dd X$ from $h_I$.

\section{Component B --- output heads, conservation map, equilibrium masking, energy and neutrinos}

\subsection{Target A (primary): signed net flux through a fixed stoichiometric layer}

Target A emits one signed scalar $\varphi_j$ per reaction and forms $\dd Y = \nu\varphi$ with $\nu$ the fixed pynucastro stoichiometric matrix. Because every column of $\nu$ satisfies $\sum_i A_i \nu_{ij} = 0$ and $\sum_i Z_i \nu_{ij} = 0$ by construction of the reaction, $\dd Y$ conserves baryon number and charge \textbf{exactly for any $\varphi$ --- including a randomly initialized, untrained network}. This is the Sturm \& Wexler result (a machine-learned replacement that ``conserves properties --- say mass, atoms, or energy --- to machine precision'' by training on fluxes and mapping through a fixed linear stoichiometric layer); accuracy only affects \emph{which} conserving update is produced, never \emph{whether} it conserves.

The dynamic range of net fluxes is handled by an \textbf{internal asinh latent}: the reaction-node latent $h_R$ is decoded to a pre-activation $z_j$, and $\varphi_j = \sinh(z_j)$ (equivalently $z_j = \asinh(\varphi_j)$). $\sinh$ is smooth and monotonic through zero and sign changes --- unlike $\log$, which is undefined for sign-changing source terms --- which is precisely D\"oppel \& Votsmeier's 2023 motivation for the asinh transform. The composition is $h_R \to z_j \to \varphi_j = \sinh(z_j) \to \dd Y = \nu\varphi$; the conserved native output is $\dd Y$, a \emph{linear} image of $\varphi$, and the nonlinearity sits upstream of the linear stoichiometric map, so conservation is untouched. The loss is computed on $\dd Y$ (with $\varphi$ as an auxiliary target). A units test on random weights must drive baryon drift $|\sum_i A_i\,\dd Y_i|$ and charge drift to $\le 10^{-12}$ per step in float64; if it does not, the $\nu$ export is wrong and training does not begin until the matrix is fixed.

\subsection{Target B (benchmark and fallback): additive null-space projection in a linear space}

Target B predicts $\dd Y$ directly and projects onto the constraint null-space with a fixed additive projection $P = I - C\Tr(CC\Tr)^{-1}C$, where $C$ is the $(A;\ Z{+}\text{electron};\ \text{lepton})$ constraint matrix (equivalently the D\"oppel--Votsmeier 2024 reduced-column-echelon null-space basis applied to nuclear invariants). The decisive lesson is that \textbf{the projection is only safe in a linear output space}: a $\log$ or signed-log output makes the sum of the inverse-transformed outputs a nonlinear function whose projection is unstable --- which is exactly the failure NuGNN documented and abandoned. That NuGNN failure is, however, \textbf{specific to its signed-log preprocessing and its inverse transform, not a universal verdict on hard conservation}; the present design's fixed-stoichiometry conservation is a structurally different mechanism and is largely orthogonal to it. Target B therefore emits $\dd Y$ linearly (asinh permitted only as an internal latent, never as the space on which the sum constraint is evaluated), and includes an explicit NuGNN-replication unit test: train a small signed-log variant and confirm the conservation-loss instability reproduces --- if it does \emph{not} reproduce, the regime differs from NuGNN's and the trap reasoning is re-examined. Target B is preferred over A only if Phase-0 shows its per-step $|\dYe| \le 2\times$ A's and it is materially simpler to train; otherwise it remains the benchmark and the fallback (and the natural target if the active-set stoichiometric matrix turns out ill-conditioned, $\mathrm{cond}(S_{\mathrm{active}}) > 10^{6}$, or if more than $\sim$90\% of reactions carry net flux below the $\Ye$ floor).

\subsection{Equilibrium masking (the load-bearing novelty)}

In QSE the forward/reverse pair of a reaction near-cancels, so the net flux is orders of magnitude below either gross rate; a mask identifies equilibrated reactions so the head does not attempt to resolve a near-zero net difference between two large learned numbers (Guidry et al.\ 2011/2013 formalize this ``microscopic equilibration''). The mechanism is a \textbf{hybrid}: a cheap, training-free Guidry partial-equilibrium test --- for a pair the analytic equilibrium abundance $\bar y_i$ is compared to the current $y_i$ and the pair is flagged when $|y_i - \bar y_i|/\bar y_i < \varepsilon$ with $\varepsilon \approx 0.01$ --- computes an imposed prior mask $m_j \in \{0,1\}$, which seeds and biases a learned \textbf{$L_0$ hard-concrete gate} (Louizos et al.\ 2018): gate logit $= \alpha(2m_j - 1) + \mathrm{MLP}(h_R)$, with the $L_0$ prior centered on $m_j$, so the network can only \emph{override} the physics prior where the data demand it and the override is penalized. The gate acts on the latent flux, $\varphi_j = g_j\cdot\sinh(z_j)$; a straight-through estimator is rejected as gradient-biased and sparsemax/$\alpha$-entmax as wrong-semantics (they normalize to a sum-to-one distribution rather than independent on/off gates).

Masking in flux space is what makes this safe: $\dd Y = \nu(g\odot\varphi)$ still satisfies $\sum_i A_i\,\dd Y_i = 0$ because the gate scales each column and the column sums vanish regardless of the scalar $g_j$ --- a structural advantage over masking in species space. Three properties bound the design. \textbf{(i)} The trajectory-varying, input-conditioned active set --- a gate whose membership re-selects per $(T, \rho, \Ye)$ state along a trajectory --- is \emph{not} inherited from the cited gating precedents, which are all static sparsification; this transfer is an unstated assumption that must be grounded in the conditional-computation / dynamic-sparsity / mixture-of-experts-routing literature and is the weakest link a referee will press. \textbf{(ii)} The two-body Guidry criterion is derived for $A{+}B \rightleftharpoons C{+}D$ pairs and does not directly capture the coupled multi-reaction (QSE-group) equilibria of Hix's group-totals (Goussis 2012 and the 2023 Michaelis--Menten criteria confirm the single-fast-timescale-per-reaction limitation); the learned gate is intended to recover that coupling, and whether it does is a Phase-0 question. \textbf{(iii)} $\varepsilon$ must be set by a Phase-0 sweep over $\{3\times10^{-3},\, 10^{-2},\, 3\times10^{-2}\}$: too tight under-masks (cancellation errors return), too loose over-masks (real slow net flux is zeroed, biasing $\Ye$), and that $\varepsilon \approx 0.01$ transfers from Guidry's $\alpha$-network tests to a pynucastro heterogeneous network is itself a hypothesis. The falsifiable fallback is to \textbf{freeze the gate to the imposed Guidry mask} (``bottleneck-only Target A'') if the learned gate regresses validation $|\dYe|$ by $>2\times$ over three seeds or if learned overrides flip $>5\%$ of memberships per step (rollout-destabilizing churn); a pure learned gate is adopted only if it beats the hybrid on $|\dYe|$ by $>25\%$ with churn $<5\%$/step.

\subsection{Lepton bookkeeping and the conservation diagnostics}

The constraint matrix is built so that strong and electromagnetic reactions conserve $Z$ and $A$ column-wise while \textbf{weak reactions change $Z$ at fixed $A$}: a charge row that includes the electron column ($Z_i$ for nuclei, $-1$ for the electron, $0$ for the neutrino) and a lepton-number row (electron $+1$, neutrino $+1$, anti-$\nu$ $-1$) are appended to the baryon row. $\Ye = \sum_i Z_i Y_i$ then evolves precisely because the weak columns have nonzero $\sum_i Z_i \nu_{ij}$ (charge moves to the leptons) while $\sum_i A_i \nu_{ij} = 0$ (baryon number fixed). The diagnostics that must be reported each step and cumulatively are: baryon drift $D_A = \sum_i A_i\,\dd Y_i$ ($\le 10^{-12}$ by construction); charge-to-lepton balance $D_Q = \sum_i Z_i\,\dd Y_i + \dd N_{\text{lepton}}$ (must close to $\le 10^{-12}$); $\sum_i X_i - 1$; and the \emph{physical} signal $\dd \Ye$ through weak reactions $= \sum_{j\in\text{weak}} \big(\sum_i Z_i \nu_{ij}\big)\varphi_j$, which must be nonzero and accurate. The first three are residuals that must vanish; the fourth is the physics that must survive --- so a design that zeroes the fourth to make the first three vanish is wrong, which is why \textbf{weak-reaction columns are excluded from the mask's eligible set} ($\beta$-decays are not in detailed-balance equilibrium in this regime); masking any weak column is, by definition, a correctness bug.

\subsection{Energy and neutrino heads}

The nuclear energy release is computed from the \textbf{same signed net fluxes that drive composition}: $\enuc = \sum_j Q_j(T)\,\varphi_j$, with the $Q$-values from pynucastro mass excesses/binding energies. Because Target A's $\varphi$ produce $\dd Y = \nu\varphi$, this flux route is algebraically identical to the composition route ($\sum_i (\text{mass excess})_i\,\nu_{ij} = -Q_j$ by definition of $Q$), so deriving $\enuc$ from the net fluxes is automatically consistent with the composition update --- a consistency the NNN lacks, where the energy and neutrino terms are separate output channels of a single shared network rather than functions of a flux it also uses for composition. (For Target B, which has no $\varphi$, only the composition route is available --- an additional reason A is preferred.) The subtle, load-bearing point is the \textbf{behavior under the mask}: a masked equilibrated pair carries near-zero \emph{net} flux $\varphi_j \approx 0$ but huge gross throughput, and because $\enuc$ uses $Q_j\varphi_j$ (net), the masked reaction contributes near-zero net energy --- which is physically correct, since an equilibrated pair releases and reabsorbs energy at near-equal rates. Masking the net flux therefore does \emph{not} corrupt $\enuc$; the danger would arise only if the head summed gross forward/reverse energies separately, which the design forbids. To keep $\enuc$ consistent across the QSE$\to$NSE transition, the $Q_j(T)$ are tabulated from the same temperature-dependent partition functions the backbone sees (not $T$-independent values), and the flux-route-versus-composition-route residual must stay $\le 1\%$ of $|\enuc|$ across the 3--4~GK band. The neutrino head is a parallel design-only channel $\varepsilon_\nu = \sum_{j\in\text{weak}} \langle E_\nu\rangle_j\,\varphi_j$ driven by the same weak net fluxes, coupling neutrino losses to the lepton bookkeeping; its tabulation-uncertainty accuracy floor is explicitly out of scope here and is the larger relative uncertainty of the two.

\section{Components C and D --- variable-\texorpdfstring{$\Delta t$}{dt} temporal head, rollout stability, and training recipe}

\subsection{Temporal head}

The baseline NNN spans $10^{-6}$--$10^{2}$~s with nine per-timestep models and runtime interpolation; the goal is to span those eight decades with \textbf{one} model. Three approaches were surveyed: $\Delta t$ as a global input feature (NuGNN's approach, simplest, composes trivially with Components A/B); a latent neural ODE on the GNN encoding (the MACE pattern --- an autoencoder plus a latent ODE solved with an ODE integrator, \emph{not} itself a graph network); and an operator head (DeepONet/FNO with $\Delta t$ as a trunk input; AMORE's mass-conserving variant uses an invertible analytical map to an $(n{-}1)$-dimensional space, structurally the same idea as Target B's null-space projection). The decision is to \textbf{ship $\Delta t$-as-a-global-input-feature on a stiffness-aware logarithmic-$\Delta t$ grid as v1} --- the lowest-risk choice, proven for variable-$\Delta t$ input by NuGNN, composing cleanly with the backbone and head --- and to design the latent neural ODE as a v2 upgrade only behind a measured stability win. The switch criterion is the whole-range falsifier: upgrade to the latent ODE if, on a constant-$(T,\rho)$ 1000-step rollout, the log-$\Delta t$-grid head's per-step $|\dYe|$ exceeds the $\approx 3\times10^{-6}$ gate at \emph{any} $\Delta t$ in $[10^{-6}, 10^{2}]$~s while a latent-ODE prototype demonstrably holds it (a high-$\Delta t$-only trigger would miss precisely the small-$\Delta t$ error-dominated drift the rollout literature warns about). The operator head is rejected if its trunk cannot be made compatible with the bipartite node structure without flattening it (which would discard the inductive bias that made NuGNN beat the FNN). The eight-decade single-model span and the 690-isotope$\to$silicon-burning transfer are labeled as transfers, not proven results.

\subsection{Conservation in the decode step}

All conservation structure lives in the \textbf{decode step, never in the latent dynamics}: the latent ODE/operator is unconstrained, and Component B's map --- the stoichiometric layer for Target A, the null-space projection for Target B --- is applied after decode, so the decoded native output is the linear conserved quantity. A latent-space soft conservation loss is rejected on the strength of MACE's verbatim finding that it ``slowed down the training immensely, because, for every pass of training data through MACE, the Jacobian of the decoder neural network needs to be calculated and multiplied by the tensor coefficients of the latent ODE, which involves many operations on large matrices.'' Any latent formulation that cannot be made exactly conserving by a fixed linear decode layer is not adopted.

\subsection{Rollout governor}

A continuous-time formulation is not a free stability win: van de Bor et al.\ 2025 showed the Branca \& Pallottini DeepONet emulator diverges under iteration through accumulated prediction error. The rollout is therefore governed by \textbf{Ono \& Sugimura's 2026 timescale-based rescaled update}, which diagnoses and fixes exactly the QSE-relevant failure --- ``if the timestep is very short, the physical change over that step can be smaller than the prediction error \ldots\ iterating such updates drives the solution away from the physical one.'' Their rule defines a characteristic timescale $\tau_\varepsilon$ at which a variable changes by a fraction $\varepsilon$ (they use $\varepsilon = 0.1$) and sets $\Delta x(\Delta t) = (\Delta t/\tau_\varepsilon)\cdot\Delta x_{\mathrm{pred}}(\tau_\varepsilon)$ when $\Delta t < \tau_\varepsilon$, else $\Delta x_{\mathrm{pred}}(\Delta t)$ --- i.e., when the requested step is error-dominated it rescales a reliable longer-step prediction down rather than trusting the short-step prediction, with $\tau_\varepsilon$ capped at the maximum learned timestep and conservation re-imposed after each update. This method comes from six-species primordial-collapse chemistry, so its adoption as the silicon-burning governor is an \textbf{out-of-regime transfer}, hedged by retaining NuGNN's cruder retry-with-smaller-$\Delta t$ heuristic as a fallback behind a measured comparison, and its commutation with Component B's conservation map (whether the rescaling preserves exactness) is an open interaction to verify.

\subsection{Rollout stability and the accumulation model}

Whether ``noise injection plus conservation defeats autoregressive divergence'' transfers from non-stiff particle dynamics to a stiff system near QSE is the core stability question, and it is symmetric: GNS random-walk noise injection, MeshGraphNets, Dynami-CAL (pairwise momentum conservation over 16\,000 rollout steps), and Brandstetter's pushforward trick are \emph{all} non-stiff particle/PDE methods, so \textbf{both} the noise-injection lever and the pushforward lever are stiff-regime hypotheses, not established remedies; Brandstetter's zero-stability ($\kappa$) result is for PDE solving and is not evidence of stiff-regime superiority. The conservation mechanism does transfer --- $A$ and $Z$ are additive linear invariants like momentum --- but near-QSE stiffness has no particle-dynamics analog: injecting noise into a near-equilibrated abundance perturbs it off the QSE manifold, where the stiff Jacobian can amplify the off-manifold component before it relaxes. The design therefore makes \textbf{pushforward/unrolled training through $K \ge 2$ steps the primary lever} (it trains on the model's own error distribution without injecting artificial off-manifold noise) and restricts GNS-style noise to channels flagged non-equilibrated by Component B's mask (a requirement imposed on the masking head), with a falsifiable rule to lock noise to non-equilibrium channels permanently if an ablation shows noise-on-all degrades the constant-$(T,\rho)$ $\Ye$ drift.

The \textbf{accumulation model is the single largest lever on the pass/fail gate}, because under systematic (linear) accumulation the per-step budget is $F/N$ ($\approx 5\times10^{-7}$ at $N = 10^{4}$ for $F = 5\times10^{-3}$) while under a zero-mean random walk it is $F/\sqrt{N}$ ($\approx 5\times10^{-5}$ at $N = 10^{4}$) --- a factor of 100 at $N = 10^{4}$. The \textbf{conservative systematic gate of $3\times10^{-6}$ (at $N \approx 1.6\times10^{3}$) remains operative} until the model is measured; the random-walk branch is not adopted as the operative gate before measurement. The measurement protocol, handed to Phase-0, is: run $M$ independent constant-$(T,\rho)$ rollouts of $N = 10^{3}$--$10^{4}$ steps recording the per-step $\Ye$ residual $r_k$; perform a cumulative-sum drift test (regress $\log|S_N|$ against $\log N$, where slope $\approx 1.0$ implies systematic bias and gate $F/N$, slope $\approx 0.5$ implies a zero-mean random walk and gate $F/\sqrt{N}$); run a zero-mean test ($t$-test/bootstrap CI on $\mathrm{mean}(r_k)$ plus the autocorrelation function, since significant positive lag-1 autocorrelation is effectively systematic even if the marginal mean $\approx 0$); and cross-check by regressing the cumulative $|\dYe|$ error against $N$. Residuals are pushed toward zero-mean by a signed/symmetric loss, an explicit bias-correction term, the exact conservation structure (Target A's stoichiometric map makes the $A,Z$ residual identically zero, isolating the $\Ye$ residual to weak reactions), and ensemble averaging --- whose effect is itself diagnostic, since it reduces a zero-mean component as $1/\sqrt{\text{members}}$ but leaves a systematic component untouched.

\subsection{Loss and uncertainty quantification}

The base loss is signed-log (asinh) MAE, augmented by a \textbf{worst-case tail term} (90th/99th-percentile or CVaR weighting), because the NNN's documented weakness was large per-isotope composition errors reported as means rather than distributions; the loss is multi-objective across composition, energy, and neutrinos and \textbf{up-weights the Fe-peak $A \approx 45$--$65$ LMP nuclei} (the electron-capture/$\beta$-decay species that set $\Ye$ --- a decrease in $\Ye$ drastically increases the iron-peak QSE abundances, so these nuclei dominate the $\Ye$ budget). The falsifiable rule is to report per-isotope error as a full distribution and to raise a nucleus's loss weight until its 99th-percentile error no longer moves $\Ye$ by more than $3\times10^{-6}$/step. Uncertainty quantification uses a \textbf{deep ensemble on the flux head} (better-calibrated tails than MC-dropout, and the spread doubles as the accumulation diagnostic above), feeding an out-of-distribution flag (ensemble-variance threshold or distance to the training manifold) that drives a \textbf{fallback-to-solver gate}: when flagged, the zone is handed back to the bbq/MESA network solver. The biggest UQ risk is the \textbf{Sobol$\to$real-MESA transfer}: error distributions measured on quasi-random Sobol compositions are not guaranteed to transfer to the thin, correlated manifold of real MESA trajectories. The mitigation is to measure error on held-out real MESA tracks (not only Sobol) and, if the real-track 99th-percentile $\Ye$ error exceeds the Sobol-measured value by more than $\approx 3\times$, to retrain with trajectory-aware sampling in the manner of DeePODE's evolutionary Monte Carlo scheme, which evolves each training sample along its local ODE trajectory.

\section{Cross-component contracts}

The three components compose through a small set of explicit contracts, stated once here rather than repeated within each component.

\begin{itemize}[leftmargin=*]
  \item \textbf{Backbone $\to$ temporal head.} The backbone must expose a graph-level (global) feature slot for $\log\Delta t$ and must preserve the bipartite isotope/reaction node structure (no flattening), so the temporal head's conditioning and any future operator trunk attach without destroying the inductive bias that distinguishes the graph from a dense baseline.
  \item \textbf{Backbone $\to$ output head.} The backbone is target-agnostic and hands off per-isotope latents $h_I$ and per-reaction latents $h_R$; Target A reads $\varphi$ from $h_R$, Target B reads $\dd X$ from $h_I$.
  \item \textbf{Output head $\leftrightarrow$ temporal head --- conservation lives in the decode.} The temporal head hands off the single \emph{linear} native quantity ($\varphi$ for A, $\dd X$ for B) and applies no nonlinear transform after the conservation map; Component B's stoichiometric map or null-space projection, and the asinh-internal nonlinearity, must live in the decode step. Conservation cannot be enforced inside an opaque latent integrator.
  \item \textbf{Output head $\to$ stability lever.} Component B's equilibrium mask is a required input to the noise-injection gate, which is restricted to non-equilibrated channels.
  \item \textbf{Rollout governor $\leftrightarrow$ conservation map.} The timescale-based rescaled update must be verified to commute with the conservation map without breaking exactness --- the one identified C-to-B interaction that could be a bug.
  \item \textbf{Mask membership $\leftrightarrow$ rollout.} A mask whose membership changes along a trajectory creates piecewise-defined dynamics; under pushforward/multi-step losses, membership flips inject discontinuities in $\dd Y$ and $\enuc$ that can destabilize the rollout gradient. The mitigations available are a membership-churn penalty, a temperature-annealed (soft$\to$hard) gate schedule synchronized with the rollout horizon, and hysteresis (two thresholds) on the Guidry mask; if churn is intrinsically high, the hybrid mask collapses to the frozen-Guidry fallback and surrenders the learned-refinement novelty.
\end{itemize}

\section{Verification status, novelty, and the load-bearing open risks}

Across three independent adversarial passes --- a backbone referee report, a citation audit centered on the output-head and masking literature, and a hostile referee audit of the temporal head and training recipe --- the citation base is unusually clean: every load-bearing reference that could be tested resolves to a real paper that broadly says what is attributed to it, including the highest-risk recent items (NuGNN, posted 3 June 2026; Padiyar et al.'s March 2026 combustion GNN; Dynami-CAL in \emph{Nature Communications}; and Ono \& Sugimura, genuinely \emph{published} as ApJ 996, 9 (2026), not a preprint mislabeled as published). No hallucinated load-bearing reference was found in any pass. The corrections the audits required were matters of attribution, labeling, internal consistency, and metadata; they are incorporated throughout the text above and are not re-enumerated here.

The \textbf{novelty} verdict is consistent and current to mid-June 2026: no published or preprinted work combines a conservation-by-construction GNN with the silicon-burning/CCSN regime, a head-to-head comparison against the Grichener 2025 baseline on identical public data, and in-the-loop MESA coupling. NuGNN establishes that the graph-on-reaction-network idea is public --- so speed of publication matters and the design leads with the two features NuGNN lacks (hard conservation by construction and the silicon-burning weak-interaction regime) --- but it is an X-ray-burst surrogate without hard conservation, the Grichener NNN is in the right regime but is a dense network without hard conservation, and the only hard-conservation reaction-network surrogate found in the adjacent literature is a combustion ANN, not a GNN and not astrophysical. RHINE (r-process heating, in-the-loop) and DeePODE (Type Ia, 3--13 species) are adjacent energy/flow surrogates, not direct competitors. Each audit notes that its forward sweep was budget-limited and that the field is moving (AMORE in October 2025, NuGNN in June 2026, conservation-loss neural emulators appearing through 2026), so the novelty check is re-run at each phase boundary rather than assumed current.

What survives verification as genuine, load-bearing uncertainty --- the risks that no literature can retire and that the design is explicitly conditional upon --- are the following. The \textbf{size-transfer hypothesis} (Component A): the GNS/MeshGraphNets precedent covers count transfer, not the new-physics neutron-rich nodes where \msn{80} fails. The \textbf{trajectory-varying mask} (Component B): the input-conditioned, re-selecting-per-state active set is not inherited from the static-sparsification precedents and needs conditional-computation grounding, and the two-body equilibrium criterion does not directly capture multi-reaction QSE-group equilibria. The \textbf{asinh conditioning} (Components B and D): $\sinh$ saturates super-exponentially for large $|z|$, and the cited motivation for asinh is sign-handling, not gradient health, so its behavior at silicon-burning dynamic range --- and whether even asinh of a genuinely tiny post-cancellation net flux loses precision --- is unproven. The \textbf{accumulation model} (Component D): the systematic-versus-random-walk question moves the per-step gate by two orders of magnitude and is the single biggest lever; if the residual is even weakly biased, the gate tightens to a level no surveyed emulator has demonstrated for a controlled scalar like $\Ye$. The \textbf{Sobol$\to$real-MESA transfer} (Component D): the largest UQ risk, since a jointly-overconfident ensemble could fail the OOD gate silently exactly where the transfer fails. The \textbf{out-of-regime rollout governor} (Component C): resting the governor entirely on a six-species primordial-chemistry method is a single point of failure, hedged but not eliminated by the retained retry fallback. And the \textbf{lepton-closure-under-error} question (Components B and D): exact $A,Z$ conservation isolates the $\Ye$ error to the weak reactions, but a biased weak-reaction flux can still drift $\Ye$ even with $A$ and $Z$ exactly conserved.

\section{Phase-0 measurement program}

The architecture is committed conditionally on a set of empirical measurements that only the bbq/MESA data can settle. These are the quantities the literature gives only as ``orders of magnitude,'' never as the distributions the design needs, and they are deferred to Phase-0 rather than asserted.

\begin{longtable}{@{}>{\raggedright\arraybackslash}p{0.37\textwidth} >{\raggedright\arraybackslash}p{0.33\textwidth} >{\raggedright\arraybackslash}p{0.22\textwidth}@{}}
\toprule
\textbf{Measurement} & \textbf{What it sets} & \textbf{Decision it gates} \\
\midrule
\endfirsthead
\toprule
\textbf{Measurement} & \textbf{What it sets} & \textbf{Decision it gates} \\
\midrule
\endhead
\midrule
\multicolumn{3}{r}{\footnotesize\itshape continued on next page}\\
\endfoot
\bottomrule
\endlastfoot
Reaction count and graph radius/diameter of \msn{80/151/204} (pynucastro export) & Depth $K \approx \lceil\text{radius}\rceil + 2$; flux-head output dimension & Component A processor depth \\
\addlinespace
Feature distributions of all proposed channels across the regime box & Normalization; dead/degenerate-channel detection & Component A feature set \\
\addlinespace
Per-channel ablation of the four added physics channels ($\ge 0.005$ MAE retention) & Which channels survive & Component A feature set \\
\addlinespace
Weight-shared vs untied processor at equal step count (within $0.01$ MAE?) & Recurrent vs untied default & Component A processor; Component C latent-ODE option \\
\addlinespace
Zero-shot \msn{80}$\to$\msn{151} transfer on neutron-rich isotopes ($\le 2\times$ internal $\Ye$ error?) & Size-transfer hypothesis & Component A size-transfer claim \\
\addlinespace
Per-reaction cancellation ratio $\kappa_r = |\mathrm{net}|/(\mathrm{gross}_f + \mathrm{gross}_r)$ distribution & How much QSE-cancellation risk the mask removes; whether Target A's reaction-space output is mostly noise & Target A vs Target B; mask design \\
\addlinespace
Condition number of the full and active-set stoichiometric matrices & Whether flux errors amplify unacceptably & Target A vs Target B (switch at $\mathrm{cond} > 10^{6}$) \\
\addlinespace
Guidry $\varepsilon$ sweep over $\{3\times10^{-3},\, 10^{-2},\, 3\times10^{-2}\}$ on silicon-burning trajectories & Mask threshold; whether $\varepsilon \approx 0.01$ transfers & Component B mask \\
\addlinespace
Mask membership churn rate per step along real $T(t), \rho(t)$ tracks & Whether the hybrid mask is rollout-compatible & Hybrid vs frozen-Guidry fallback \\
\addlinespace
Empirical baryon/charge drift of Target B's projection at realistic dynamic range & Target B stability under 20-decade abundance spread & Target B viability \\
\addlinespace
$\enuc$ flux-route $-$ composition-route residual across the 3--4~GK QSE$\to$NSE band ($\le 1\%$?) & Partition-function consistency & Component B energy head \\
\addlinespace
$\Ye$-residual accumulation model: slope of $\log|\text{cumulative}|$ vs $\log N$ ($\approx 1$ systematic, $\approx 0.5$ random-walk) & Per-step gate ($3\times10^{-6}$ vs $5\times10^{-7}$ vs $5\times10^{-5}$) & The pass/fail gate --- the single biggest lever \\
\addlinespace
Sobol$\to$real-MESA error-transfer ratio (99th-percentile $\Ye$ error, real $\div$ Sobol) & Whether trajectory-aware sampling is mandatory & Component D sampling strategy (retrain at $>3\times$) \\
\addlinespace
Per-isotope error \emph{distribution} for Fe-peak $A \approx 45$--$65$ nuclei & Worst-case isotope's effect on $\Ye$ & Component D loss weighting \\
\addlinespace
Whether the timescale-based update holds the gate at $\Delta t \ge 0.1$~s; whether noise-on-non-eq beats noise-on-all; whether the log-$\Delta t$-grid head needs the latent-NODE upgrade & Rollout governor, noise gate, temporal-head v1/v2 & Components C/D \\
\end{longtable}

Several items also require local code-level confirmation before sizing or training: the exact \msn{80}/\allowbreak\msn{151} isotope lists and which $\Ye$-controllers (especially the neutron-rich $\beta$-decay partners) each contains; which weak-rate tables bbq/MESA r23.05.1 actually loads and their off-grid-edge extrapolation behavior; MESA's average-neutrino-energy handling for the neutrino head; whether MESA's softwired inverse rates reproduce equilibrium cleanly enough that $\kappa_r \to 0$ without a spurious floor (a screen that must precede any kill-test conclusion); and the precise meaning of the Farmer 2016 ``30\%/10\%'' figures.

\section{Conclusion}

The three architecture runs jointly specify a conservation-by-construction GNN emulator for silicon burning that is implementable end-to-end and that keeps its central differentiator --- hard baryon and charge conservation --- without walking into the QSE catastrophic-cancellation regime that defines the problem. The backbone is a heterogeneous encode--process--decode network on the bipartite isotope/reaction graph, with the four physics channels the $\Ye$ regime requires, untied message-passing blocks as the proven default, depth gated on validation MAE rather than a confounded over-smoothing alarm, and size transfer carried by a continuous $(Z,N)$ embedding whose extension to chemically new isotopes is honestly flagged as the backbone's biggest bet. The output head predicts a signed net flux mapped through a fixed stoichiometric matrix, so conservation holds to machine precision for any prediction while $\Ye$ evolves correctly through the weak columns of the constraint matrix; a hybrid equilibrium mask --- a cheap physics prior refined by a learned gate, acting on the latent flux and never on weak reactions --- keeps the head from resolving the near-zero net flux of an equilibrated pair, and the energy release, built from the same net fluxes, is consistent with the composition update by construction. The temporal head ships a single variable-$\Delta t$ model on a logarithmic grid, keeps all conservation in the decode, governs the rollout with a timescale-based update that refuses to trust error-dominated short steps, trains with pushforward losses and a worst-case tail term that targets the Fe-peak weak-rate nuclei, and gates its own out-of-distribution failures back to the solver.

The design has been verified three times over and emerges with its citation base intact, its novelty confirmed for the silicon-burning conservation niche as of mid-June 2026, and a short, well-posed list of load-bearing uncertainties that no literature can resolve: whether size transfer holds across new weak physics, whether the per-step $\Ye$ residual accumulates systematically or as a random walk, whether Sobol-measured error transfers to real stellar trajectories, whether the trajectory-varying mask is stable and properly grounded, and whether the asinh latent and the out-of-regime rollout governor behave at silicon-burning dynamic range. Each of these resolves into a specific Phase-0 measurement with a stated threshold and a stated alternative --- switch the target if the cancellation ratio or the condition number says so, freeze the mask if its membership churns, upgrade to a latent ODE if the grid head drifts, tighten or relax the gate as the accumulation slope dictates, retrain on trajectory-aware samples if Sobol fails to transfer --- so the architecture is committed conditionally, with explicit instructions to switch rather than persist, and with the operative accuracy gate held at the conservative $|\dYe| \lesssim 3\times10^{-6}$ until the measurements say otherwise.

\section*{Key references (post-verification)}
\addcontentsline{toc}{section}{Key references}

\paragraph{Baselines and the closest competitor.} Grichener et al.\ 2025, \emph{ApJS} 279, 49 (arXiv:2503.00115) --- the NNN baseline, a single shared dense network, log-softmax ($\sum X = 1$) composition head, nine per-timestep models, $\Ye$ improvement 280--660\% and energy 250--750\% over a 22-isotope net. Kim, Chae, Ko, Mumpower \& Smith 2026, arXiv:2606.04491 (preprint) --- NuGNN, the first bipartite isotope/reaction GNN (690-isotope XRB), signed-log $\dd X$, GAT+sigmoid-gate attention, flux preprocessor, no hard conservation (abandoned conservation loss/activation as log-space-specific instability), retry-with-smaller-$\Delta t$.

\paragraph{Architecture / encode--process--decode.} Battaglia et al.\ 2018 (arXiv:1806.01261); Sanchez-Gonzalez et al.\ 2020, GNS (arXiv:2002.09405, ICML); Pfaff et al.\ 2020, MeshGraphNets (arXiv:2010.03409, ICLR 2021); Veli\v{c}kovi\'c et al.\ 2018, GAT; Brody et al.\ 2021, GATv2 (arXiv:2105.14491); over-smoothing --- Li, Han \& Wu 2018 (AAAI); Oono \& Suzuki 2020 (ICLR); Rusch et al.\ 2023 survey (arXiv:2303.10993). E$(n)$/SE(3) rejected: Satorras et al.\ 2021 (arXiv:2102.09844); Fuchs et al.\ 2020. Chart-of-nuclides static-property GNN: Choi, Mitra et al., \emph{Phys.\ Rev.\ C} (2025) (arXiv:2404.02332). Hypergraph alternatives rejected: ChemHGNN (arXiv:2506.11041); Rxn-Hypergraph (arXiv:2201.01196).

\paragraph{Conservation layers and output heads.} Sturm \& Wexler 2020, \emph{GMD} 13:4435--4442 (flux target + fixed linear layer $\Rightarrow$ machine-precision conservation) and 2022, \emph{GMD} 15:3417--3431; D\"oppel \& Votsmeier 2023, \emph{React.\ Chem.\ Eng.} 8:2620 (internal asinh latent) and 2024, \emph{Proc.\ Combust.\ Inst.} 40:105507 (null-space atom-balance layer); Ji \& Deng 2021, CRNN (\emph{J.\ Phys.\ Chem.\ A} 125:1082; arXiv:2002.09062); Kircher \& Votsmeier 2025, \emph{J.\ Phys.\ Chem.\ Lett.} 16:4715 (positivity-preserving projection); Fan et al.\ 2022 (arXiv:2207.10628; \emph{ApJ} 940:134; energy-from-composition identity).

\paragraph{QSE physics and equilibrium masking.} Hix \& Thielemann 1996, \emph{ApJ} 460:869 (Silicon Burning I; QSE structure, ${}^{45}\mathrm{Sc}(p,\gamma){}^{46}\mathrm{Ti}$ bottleneck) and 1999, \emph{ApJ} 511:862 (Silicon Burning II; QSE for $T \gtrsim 3$~GK); Hix et al.\ 2007, \emph{ApJ} 667:476 (QSE-reduced network); Guidry et al.\ 2011/2013, \emph{Comput.\ Sci.\ Disc.} 6:015003 (partial-equilibrium criterion); Goussis 2012, \emph{Combust.\ Theor.\ Model.} 16:869, and the 2023 Michaelis--Menten criteria paper (single-fast-timescale limitation); gating primitives --- Louizos et al.\ 2018, $L_0$ hard-concrete (ICLR; arXiv:1712.01312); Bengio et al.\ 2013, STE; Martins \& Astudillo 2016 / Peters et al.\ 2019, sparsemax/entmax; Heger, Woosley, Mart\'inez-Pinedo \& Langanke 2001, \emph{ApJ} 560:307 (final $\Ye$ set during Si shell burning).

\paragraph{Temporal head, rollout stability, and training.} Maes et al.\ 2024, MACE, \emph{ApJ} 969:79 (arXiv:2405.03274; autoencoder + latent ODE, not a GNN; decoder-Jacobian warning; mean speedup $\times 26$); Goswami et al.\ 2023, \emph{CMAME} 419:116674 (DeepONet/PoU-DeepONet); Nath et al., AMORE (arXiv:2510.12999); van de Bor et al.\ 2025, \emph{OJAp} (arXiv:2503.10736; DeepONet instability under iteration); Branca \& Pallottini 2024, \emph{A\&A} 684, A203; Ono \& Sugimura 2026, \emph{ApJ} 996, 9 (arXiv:2508.16114; timescale-based rescaled update); Brandstetter et al.\ 2022 (arXiv:2202.03376, ICLR; pushforward trick); Sharma \& Fink 2026, Dynami-CAL, \emph{Nat.\ Commun.} 17:1045; Holdship et al.\ 2021, \emph{A\&A} 653, A76; Sulzer \& Buck 2023 (arXiv:2312.06015); Yao et al.\ 2025, DeePODE, \emph{ApJ} (arXiv:2504.14180; evolutionary Monte Carlo trajectory sampling).

\paragraph{Infrastructure and physics base.} pynucastro --- Willcox \& Zingale 2018, \emph{JOSS} 3:588, and Smith et al.\ 2023, \emph{ApJ} 947:65 (NetworkX export, REACLIB fits, weak-rate tables in $(T, \rho\Ye)$, detailed-balance reverse rates, partition functions, screening); Farmer et al.\ 2016, \emph{ApJS} 227:22 ($\approx$127-isotope convergence floor); Rauscher \& Thielemann 2000 (detailed-balance reverse rates).

\end{document}
